% Fact check of the CBC Daybreak interview of October 9, 2026 with
% councillor Alexandre Teodoresco on rue de Terrebonne. Compile with
% pdflatex (twice), or run build.py.
\documentclass[11pt,letterpaper]{article}
\usepackage[T1]{fontenc}
\usepackage[utf8]{inputenc}
\usepackage[scaled=0.95]{helvet}
\renewcommand{\familydefault}{\sfdefault}
\usepackage[letterpaper,top=0.9in,bottom=0.95in,left=1in,right=1in]{geometry}
\usepackage{microtype}
\usepackage{enumitem}
\usepackage{titlesec}
\usepackage{fancyhdr}
\usepackage{xcolor}
\usepackage{booktabs}
\usepackage{longtable}
\usepackage{array}
\usepackage{textcomp}
\usepackage[hidelinks]{hyperref}
\urlstyle{same}
\newcommand{\docversion}{2026-10-16-17-02-24} % set by build.py (America/Toronto)
\hypersetup{pdfsubject={version \docversion}}

\setlength{\parindent}{0pt}
\setlength{\parskip}{0.5em}
\linespread{1.03}

\titleformat{\section}{\normalfont\large\bfseries}{\thesection}{0.6em}{}[\vspace{1pt}\titlerule]
\titlespacing*{\section}{0pt}{1.1em}{0.4em}
\titleformat{\subsection}{\normalfont\bfseries}{\thesubsection}{0.6em}{}
\titlespacing*{\subsection}{0pt}{0.8em}{0.25em}

\pagestyle{fancy}
\fancyhf{}
\renewcommand{\headrulewidth}{0pt}
\fancyfoot[C]{\footnotesize\color{gray}CBC Daybreak, October 9, 2026 -- fact check -- page \thepage}
\fancyfoot[R]{\tiny\color{gray}v.\,\docversion}

% ---- Links ----
\newcommand{\repourl}{https://github.com/OpenMobilityData/Terrebonne-open-sources}
\newcommand{\repo}[1]{\href{\repourl}{#1}}
\newcommand{\doc}[2]{\href{\repourl/blob/main/#1}{#2}}
\newcommand{\arurl}{https://github.com/OpenMobilityData/Terrebonne-access-requests}
\newcommand{\pagesurl}{https://openmobilitydata.github.io/Terrebonne-access-requests}
\newcommand{\cliptime}[1]{\textsf{\small[#1]}}
% \pg{pages}: page of bundle 01 of access request 2026-581, linked
\newcommand{\pgl}[2]{\href{\pagesurl/2026-581/raw/01_emails_mandate_counts.pdf\#page=#1}{\textsf{\small[2026-581, p.~#2]}}}
\newcommand{\pg}[1]{\pgl{#1}{#1}}
% \exp{page}: page of the 2023 EXP study (open-sources repository)
\newcommand{\exppg}[1]{\href{\repourl/blob/main/2023-exp-mobility-study/raw/EXP_Etude_de_mobilite_rue_de_Terrebonne_2023.pdf}{\textsf{\small[EXP 2023, p.~#1]}}}
\newcommand{\quot}[1]{\begin{quote}\small #1\end{quote}}
% Statements from the October 6 question period: machine transcript of the
% webcast, to be checked against the video before use. Set \qatrue to include.
\newif\ifqa\qafalse

\begin{document}

\begin{center}
{\Large\bfseries CBC \emph{Daybreak}, October 9, 2026: the councillor's statements checked against the record}\\[0.4em]
{\large Fact check of the interview on rue de Terrebonne}\\[0.6em]
Richard Hoge, Côte-des-Neiges--Notre-Dame-de-Grâce -- October 16, 2026\\
Published with the transcript at \repo{\repourl}
\end{center}

\section*{Summary}

On October~9, 2026, CBC Radio's \emph{Daybreak} (Montréal) interviewed
Alexandre Teodoresco, city councillor for the Loyola district of
Côte-des-Neiges--Notre-Dame-de-Grâce, about the borough's plan to convert
one direction of the protected bike lane on rue de Terrebonne to parking in
winter. Each statement of fact in the interview was compared with the
public record, and the points below say what the record shows, including
what a statement leaves out: the 508 pages released by the borough under access request
2026-581, the borough's counter data released under requests 2025-697 and
2026-300, the City of Montréal's open data, the 2023 EXP mobility study,
council records, and press coverage. Times in square brackets are minutes
and seconds from the start of the CBC clip; the full transcript is
published alongside this document and has been checked against the audio.
This document compares
statements with records, and with what the councillor said at the borough
council on October~6, three days earlier; it does not assess intent.

\begin{enumerate}[nosep]
  \item \begin{minipage}[t]{\linewidth}\textbf{``Less than one an hour'' \cliptime{03:11} is the smallest reading the counter allows, and a rate is the wrong measure.} The borough's counter at Kensington recorded 4{,}891 bicycle passages between December~1 and March~31: 40.4 a day, 1.68 an hour over 24 hours, both directions. The figure holds only for one direction averaged over the whole day and night (eastbound 0.96, westbound 0.73). On weekdays in the City's winter period for bike lanes, November~16 to March~31, the counter averaged 7.9 an hour at 08:00 and 6.3 at 16:00, and no daytime hour fell below 2.3. What bears on safety is not the rate but the number of trips exposed: 4{,}891 over the four months, 2{,}783 of them eastbound.\end{minipage}\vspace{0.35em}
  \item \begin{minipage}[t]{\linewidth}\textbf{The ``8'' \cliptime{07:03} is a six-hour sum at one intersection on one day under a winter weather advisory, and the same slide carried higher figures.} His explanation, ``the sum of all the hours during the rush hour… six hours,'' matches the slide \pg{62}: Grand, Terrebonne axis, January~29, 2025, 06:30--09:30 and 15:30--18:30 \pg{34}. Not said: the count is one day's, taken while an Environment Canada advisory urged residents to postpone non-essential travel; the slide gives 10 at Cavendish; the twelve-hour count at Grand was 12 \pg{85}; and the borough's own counter, whose table went to the consultant on August~28 \pg{19}, gives 40 a day for the same months. On October~6 the figure was ``8 passage… 4 cyclists.''\end{minipage}\vspace{0.35em}
  \item \begin{minipage}[t]{\linewidth}\textbf{The source he named as ``the source of truth'' \cliptime{07:19} does not exist on the open-data portal.} The City's \emph{Compteurs cyclistes permanents} dataset lists 128 sites, 77 of them Eco-Counter; none is on Terrebonne. The only Terrebonne site is a City detector at Girouard whose one Terrebonne channel counts westbound traffic, from October~23, 2025. The borough's Eco-Counter loop at Kensington is not on the portal, and the borough had not mentioned it in public before a resident obtained its record under access request 2025-697 (answered December~19, 2025); its data have reached the public only through that request and requests 2026-300 and 2026-581.\end{minipage}\vspace{0.35em}
  \item \begin{minipage}[t]{\linewidth}\textbf{``From 40 a day to 30 a day during December to March'' \cliptime{07:34} matches the borough's own Kensington counter, not the open-data site he named, and leaves out March.} Kensington, both directions: December 40.0, January 32.7, February 34.9 a day; March 53.5, above the range he gave and inside the season he proposes to convert. The Girouard detector, westbound only: December 32.5, January 22.6, February 29.1, March 24.7.\end{minipage}\vspace{0.35em}
  \item \begin{minipage}[t]{\linewidth}\textbf{``We have the data for many years'' \cliptime{03:33} is one continuous winter and single-day counts.} Continuous counting on Terrebonne began on July~26, 2025 (open-data detector: October~23, 2025). Before that: camera counts at intersections on single days (Miovision, November 2024 and January~29, 2025), EXP's one day of counting on October~4, 2022 in the two peak periods \exppg{13}, and Telraam counts to December 2024 \pg{65}. The borough's own deck averages five counts for motor vehicles \pg{60} and relies on one day for bicycles, the day under the advisory. On October~6: ``data for 3 years in a row.''\end{minipage}\vspace{0.35em}
  \item \begin{minipage}[t]{\linewidth}\textbf{``About 250 parking spots'' removed \cliptime{01:08} is close to the record; the study dates from 2023, the lane from 2024.} The 2023 EXP study counts 475 spaces on Terrebonne, 375 once the 5-metre clearances are applied, and 171 in the retained design \exppg{100}: a reduction of 304, or 204 against the 375. No as-built count is in the record. In 2020 a press report gave 250 spaces removed by that year's temporary lanes (\emph{Métro}, September~1, 2020).\end{minipage}\vspace{0.35em}
  \item \begin{minipage}[t]{\linewidth}\textbf{``At least 50 to 60 parking spots'' \cliptime{01:28} is in no document, rests on no measurement, and sits oddly beside the entire stretch he confirmed.} The figure appears nowhere in the 508 pages of request 2026-581, which contain no inventory, no occupancy count and no estimate of spaces created; by his own account it is ``a bit too early'' and ``a few rounds'' with the services remain. Asked whether the whole stretch is concerned, he said yes \cliptime{09:01}. On EXP's figures the side to be converted held about 190 spaces after clearances; 50 to 60 over 2.45~km is about one space for every 40 to 50~m of lane removed. The only occupancy survey on record is EXP's, two days in autumn 2022, in the evening: 35--70\% on weekdays \exppg{54}.\end{minipage}\vspace{0.35em}
  \item \begin{minipage}[t]{\linewidth}\textbf{The route he described \cliptime{08:19} is the 2024 route, and it places the City's withheld analysis on the path.} ``We vote here at the borough, then a vote at the Conseil Municipal, then it comes back to the borough'' is how the lane was approved: CA24 170018 (February~5, 2024), CM24 0179 (February~19), CA24 170131 (June~3). City council's December sitting is on the 14th and 15th; no borough or City agenda through October~19 carries the item. The central-city documents in the file, the City's June~17 analysis (``plusieurs compromis importants… sécurité'' \pg{6}) and the fire department's opinion \pg{5}, both withheld, are not mentioned in the interview. The 2024 marking was scheduled for June; this change is for December.\end{minipage}\vspace{0.35em}
  \item \begin{minipage}[t]{\linewidth}\textbf{``So for now, we're looking at it in a way that will make it permanent'' \cliptime{12:22} is a third description of the measure, and the word is ambiguous.} The service request of August~20 calls it ``mesure transitoire pour l'hiver 2026-2027'' \pg{4}; on October~6 it was ``a test''; the August~13 mandate was a year-round layout, a two-way lane on one side with parking restored \pg{1}. The host's question concerned the winter change; the answer can mean a removal every winter or the year-round layout, and the interview does not say which. His stated trigger for revisiting it: winter bicycle traffic doubling.\end{minipage}\vspace{0.35em}
  \item \begin{minipage}[t]{\linewidth}\textbf{``The final study… still need to be submitted to us'' \cliptime{00:38} came three days after the announcement and a week after the note was due.} EXP's final note was due on October~2 \pg{12}; a draft was transmitted on September~17 \pg{167}; the outcome was announced on October~6; the note is withheld as part of a decision process still under way (s.~39). A request for the final note is pending (answer due November~7).\end{minipage}\vspace{0.35em}
  \item \begin{minipage}[t]{\linewidth}\textbf{``Reduce vehicle traffic by two-thirds'' \cliptime{09:56} is right; what remains would share one lane between parked cars.} 6{,}300 vehicles a day in 2022 between Cavendish and Grand \exppg{41}; 1{,}375 a day on the Kensington counter's first year ($-78$\,\%). Under the winter layout a cyclist in the converted direction rides in the single driving lane, 3.0 to 4.0~m on the design plans annexed to the 2023 study, between parked cars on both sides; the lane is as wide as each of the two lanes the street had before 2024, but there is no longer a second lane or a centre line to cross to give the one metre the \emph{Code de la sécurité routière} requires when passing. The driving lane's width under the winter layout is not in the release.\end{minipage}\vspace{0.35em}
  \item \begin{minipage}[t]{\linewidth}\textbf{Compared with October~6.} In his address to the borough council three days earlier, the councillor described the data as ``the official data provided by the borough'' (on \emph{Daybreak}: ``the open data website''); the winter change as a 99\,\% drop, ``from 700 during the good months to 8'' (on \emph{Daybreak}: 30 to 40 a day, and the 8 as a six-hour sum); ``data for 3 years in a row'' (``many years''); ``more than 30 meetings'' (``more than 40''); and the measure as ``a test'' (``for now… permanent''). No number of parking spaces was given in the address. The statements are set side by side in section~3.\end{minipage}\vspace{0.35em}
\end{enumerate}

\section{What was checked, and how}

\textbf{The interview.} CBC Radio One, \emph{Daybreak} (Montréal), clip
``Terrebonne bike path in CDN to shut down in winter'', dated October~9,
2026, 13~min~32~s. The transcript published with this document was made by
machine transcription and then checked word by word against the audio.

\textbf{Sources.} Access request 2026-581 (borough response of October~9,
2026; bundle~01, 167~pp., cited by page); the borough's Eco-Counter data at
Terrebonne/Kensington (daily totals, 2026-581 pp.~74--81; hourly by
direction, request 2026-300); the City's open data (\emph{Compteurs
cyclistes permanents}, site list and 2025--2026 files, and
\emph{Comptages des véhicules, cyclistes et piétons}, downloaded October~16,
2026); the 2023 EXP mobility study and the 2024 council records in this
repository; the borough's budget document of October~6, 2026; the
councillor's October~6 address to the borough council; press reports.

\textbf{Terms.} In section~2, \emph{consistent} means that the record
supports the statement as made, and \emph{cannot be verified} that the
record does not bear on it. A statement can be accurate in its terms and
still leave out what the record also shows; the summary points say what
was left out.

\section{Findings in detail}

\subsection{Cycling counts}

\textbf{Winter rate \cliptime{03:11}.} Asked by the host whether ``less
than one cyclist per hour is using the bike path on either side'', the
councillor answered ``Indeed.'' Kensington counter, December~1, 2025 --
March~31, 2026: 4{,}891 passages (2{,}783 eastbound, 2{,}108 westbound) in
121 days. Per hour over 24 hours: 1.68 both directions, 0.96 eastbound,
0.73 westbound. Weekdays 07:00--19:00 (87 days, 1{,}044 hours): 3{,}318
passages, 3.18 an hour. By month, both directions: December 40.0, January
32.7, February 34.9, March 53.5 a day. The westbound channel undercounts
from about March~1, 2026 (see the 2026-300 note in the
\href{\arurl}{access-requests repository}); the westbound figures are lower
bounds. City detector at Girouard, westbound only: 3{,}150 passages in 115
days with data, 27.4 a day, 1.14 an hour. Weekdays November~16 --
March~31, both directions, by hour of day: 7.9 at 08:00, 6.3 at 16:00,
6.0 at 17:00; the lowest daytime hour, 11:00, 2.3; the 24-hour mean 2.6.

\textbf{The ``8'' \cliptime{07:03}.} Slide ``Nombre de vélo aux heures de
pointes axes est/ouest'' \pg{62}: Terrebonne/Grand, January 2025, 8
($-95$\,\% against October 2022); Terrebonne/Cavendish, 10. Peak hours as
defined in the borough's presentation \pg{34}: 06:30--09:30 and
15:30--18:30. The underlying Miovision table for Grand, January~29, 2025, is
at \pg{85}. The same day's twelve-hour count at Grand on the Terrebonne axis
was 12. October~6 address: ``We go from 7700 [700] during the good months
to 8. That's 8 passage, right? So that's 4 cyclists.'' The day's Miovision
tables record every mode: over the same six hours, 617 cars on the two
Terrebonne approaches and 605 pedestrians, against 8 bicycles; set against
November~20, 2024, cars were down 23\,\%, pedestrians 26\,\% and bicycles
95\,\% \pg{82}--\pg{88}, \pg{131}--\pg{137}. Environment and Climate Change
Canada had issued a winter weather advisory for Montréal that morning.

\textbf{The open data source \cliptime{07:19}, \cliptime{03:36}.} Site list
of \emph{Compteurs cyclistes permanents}, October~16, 2026: 128 sites (77
Eco-Counter, 51 City). On Terrebonne: \texttt{det-00118-01}, \texttt{-02},
\texttt{-03}, ``Girouard/Terrebonne'', source ``Ville''; only
\texttt{det-00118-01} records a Terrebonne direction (``Ouest''). Daily data
start October~23, 2025. The intersection-count dataset
(\emph{Comptages des véhicules, cyclistes et piétons}) has Terrebonne counts
up to September 2023 and none after. The borough's Kensington loop is the
Eco-Counter referred to in the borough engineering division's email of
June~10, 2026 (``Nous disposons également d'un compteur permanent
Éco-Compteur sur le corridor'' \pg{8}); its data are not on the portal.
The author has found no public borough document or statement mentioning
the counter before its record was obtained under access request 2025-697
(filed November~27, 2025; answered December~19, 2025), which also released
the equipment documentation.

\textbf{Years of data \cliptime{03:33}.} On October~6 the councillor
referred to ``having data for 3 years in a row''. Records available: EXP
counts of 2022; Miovision intersection counts from July 2024 to September
2025 \pg{82}--\pg{164}; Telraam counts to December 2024 \pg{65}; Kensington
loop from July~26, 2025; open-data detector from October~23, 2025. EXP's
2022 counts were made on one day, Tuesday October~4, 2022, in the two peak
periods, at 17 intersections \exppg{13}. The borough's deck averages five
counts for motor vehicles \pg{60} and relies on one day for bicycles.

\textbf{Children in winter \cliptime{09:22}.} ``I've very, very, very
rarely seen any kids use it during winter.'' No counter in the record
records age. Cannot be verified.

\subsection{Parking}

\textbf{About 250 spaces removed \cliptime{01:08}.} EXP 2023: 230 spaces
westbound and 248 eastbound on Terrebonne \exppg{46}; 475 in the current
situation, 375 after the 5-metre clearance rule, 171 in scenario~3, the one
retained \exppg{100}. Project approval: CA24 170018 (February~5, 2024) and
CM24 0179 (February~19, 2024); deployment May--June 2024 (borough file
1246880002). \emph{Métro}, September~1, 2020, on the temporary lanes of that
year: they ``a éliminé 250 places de stationnement des deux côtés de la
rue''.

\textbf{50 to 60 spaces \cliptime{01:28}.} Not in the 508 pages. The
service request asks EXP to identify the spaces the layout ``pourrait
permettre de recréer'' \pg{4}; EXP's result is withheld. Item C4 of the
request, which asked for parking inventories and occupancy surveys, was
answered with traffic counting tables. For comparison only: on EXP's
figures, the difference between 375 and 171 spaces is 204, about 100 to a
side; the record does not say which blocks the conversion would cover. On
\emph{Daybreak} the councillor confirmed that the whole stretch is
concerned \cliptime{09:01} and said the spaces would be ``reinstated''
where they can be \cliptime{08:46}; at 6 to 6.5~m a space, 50 to 60
spaces occupy 300 to 400~m of a 2.45~km lane.

\textbf{Occupancy \cliptime{05:41}.} The councillor used ``only 60\% of the
parking used'' as an example of a figure he would not rely on. EXP surveyed
Thursday September~29 and Saturday October~1, 2022, in three one-hour
periods ending by 19:30 \exppg{13}: ``le taux d'occupation varie entre 35\%
et 70\% la semaine et entre 10\% et 60\% la fin de semaine'' \exppg{54}.
No later, winter or overnight measurement is in the record.

\subsection{Process and timing}

\textbf{Three votes \cliptime{08:19}.} Charter of Ville de Montréal,
art.~85: a borough offer to provide a service within the City council's
jurisdiction takes effect when the City council accepts it. Terrebonne is
on the arterial network (2.54~km marked arterial in the City's RAAV layer,
covering the lane). 2024 sequence: CA24 170018 (offer), CM24 0179
(acceptance), CA24 170131 (June~3, 2024, ordinance OCA24 17042, one-way).
Agendas checked: borough, June~1, July~6, September~8 and October~6, 2026
(regular) and October~6 (budget); City council, August~24, September~21
and October~19, 2026. None has a decision item on Terrebonne. The
councillor's statement card says the proposal requires consideration by
``the borough council and City Council''. The interview does not mention
the City's analysis of June~17 or the fire department's opinion of
July~13, the two central-city documents in the file, both withheld
\pg{5}, \pg{6}. City council's December sitting is on the 14th and 15th.
The 2024 signage and marking were scheduled for late June (borough file
1246880006); the change announced is to be made before Christmas.

\textbf{Permanent \cliptime{12:22}.} Full answer to the question ``what are
the chances the borough makes this change permanent?'': ``We certainly don't
want to be changing, you know, the layout of all the bike paths every year…
But, yeah, so for now, we're looking at it in a way that will make it
permanent. But if there's really a big change in the data… we'll certainly
look at it.'' October~6: ``So this is what we're going to propose. It's
going to be a test.'' August~13 mandate from the borough mayor's office: a
``piste cyclable bidirectionnelle d'un côté'' with ``le retour du
stationnement d'un côté de la rue'' \pg{1}. August~20 service request:
``Volet 1 -- existant… (mesure transitoire pour l'hiver 2026-2027)'' and
``Volet 2 -- faisabilité d'une piste cyclable bidirectionnelle
unilatérale'' \pg{4}. The host's question concerned the winter change; the
answer does not say whether ``permanent'' means a removal every winter or
the year-round layout of the August~13 mandate.

\textbf{The final study \cliptime{00:38}.} EXP's offer sets the final note
for October~2, 2026 \pg{12}; the kick-off form lists ``NT finale
2-Oct-26'' \pg{26}; the draft was transmitted September~17 \pg{167}. The
response of October~9 withholds the note under ss.~9, 37 and 39 of the
Access Act. In the October~6 question period (machine transcript, not yet
checked) the councillor referred to having received ``le dernier rapport''
the week before. The note was thus a week overdue on the day of the
interview, and the outcome it was commissioned to assess had been
announced three days earlier.

\textbf{Reopening date \cliptime{13:08}.} ``I look at the data and the
weather.'' No reopening date appears in the release; the councillor
described the season as ``December to March'' \cliptime{07:34}.

\textbf{Which side \cliptime{00:38}.} The orientation given is to keep the
lane that runs against motor traffic. The service request asks for the
maintenance of one lane and the removal of the other without naming a side
\pg{4}; the kick-off form asks EXP to consider the direction of motor
traffic against the direction of cyclists \pg{26}. No safety analysis of
this option is in the release; the only safety conclusion that survives is
the borough's summary of the City's June~17 analysis of earlier concepts,
``plusieurs compromis importants sur le plan de la sécurité'' \pg{6}.

\subsection{Traffic and street type}

\textbf{Two-thirds less traffic \cliptime{09:56}.} EXP 2023, 2022 daily
traffic between Cavendish and Grand: 1{,}800 westbound and 4{,}500
eastbound, 6{,}300 in all \exppg{41}. Kensington counter, first year
(July~26, 2025 -- July~25, 2026): 501{,}817 motor vehicles, 1{,}375 a day
($-78$\,\%). Peak-hour entries at Grand: 2{,}101 in October 2022, then
770--1{,}138 ($-46$ to $-63$\,\%) \pg{34}.

\textbf{The lane that remains.} The design plans annexed to the 2023 study
give the one-way driving lane 3.0 to 4.0~m, mostly 3.5 to 3.6~m
\exppg{346}, beside a 2.3 to 2.5~m parking lane and the protected lanes;
before 2024 the street was two-way, one lane each way, with parking on both
sides (borough file 1246880006: ``la circulation véhiculaire est permise
dans les 2 directions avec 1 bande de stationnement de chaque côté''). Under
the winter layout the converted side becomes parking, so a cyclist in that
direction rides in the single driving lane between parked cars on both
sides, and a driver passing must give one metre (\emph{Code de la sécurité
routière}, art.~341) within that lane, there being no second lane or centre
line to cross. The width of the driving lane under the winter layout is not
in the release; the June concepts described to the City had a traffic lane
``très étroite (moins de 4 mètres)'' between two parking rows \pg{8}, a
different layout.

\textbf{``Les rues vélo'' \cliptime{10:14}.} A \emph{vélorue} is designated
by municipal by-law under the \emph{Code de la sécurité routière}
(arts.~496.1--496.8; 30~km/h). Québec's guidance: at most 1{,}000 vehicles
a day, ideally 500, not on the trucking network. Montréal's design guide:
under 1{,}000 vehicles a day, local trucking only, 30~km/h. Terrebonne:
1{,}375 vehicles a day in the counter's first year; arterial network;
trucks prohibited except local delivery; 30~km/h. Terrebonne is not on the
City's list of vélorues.

\subsection{Other statements}

\textbf{Executive committee \cliptime{00:25}.} The host's description
(``responsible for optimization in municipal performance and innovation'')
matches the City's announcement of November~18, 2025. Consistent.

\textbf{Since 2020 \cliptime{00:08}.} The host's ``first put in place in
2020'': temporary lanes were installed in 2020 and removed in September
2020 after a council motion; the present lanes date from 2024. Partly
consistent.

\textbf{Meetings \cliptime{02:37}.} ``More than 40 meetings and work
sessions''; on October~6, ``more than 30 meetings on this topic with the
services''. The release names a meeting with the borough mayor on June~18
\pg{8}, the kick-off on September~3 \pg{26} and a working meeting on
September~11 \pg{167}; for agendas or minutes of follow-up meetings the
response states that no documents were found. Cannot be verified.

\textbf{School zones \cliptime{10:52}.} ``\$1.5 million over the next two
years.'' The borough's budget document of October~6, 2026 lists
``760~000~\$ spécifiquement pour sécuriser et apaiser les abords des
écoles'' and an addition of 700{,}000~\$ over two years to the budget for
traffic-calming measures. The figure of \$1.5 million does not appear.
Cannot be verified as stated.

\textbf{Campaign \cliptime{01:53}.} ``4,000 doors'': no source. That the
lane was a campaign issue is consistent with press coverage (\emph{Montreal
Gazette}, November~26, 2025 and October~7, 2026).

\section{Compared with the October~6 council meeting}

On October~6, 2026 the councillor addressed the borough council on
Terrebonne and answered questions from the public. The address is quoted
from the webcast captions (``Taliban'' corrected to Terrebonne; the
captions' ``7700'' is 700, as the slides show). His statement card of the
same day is quoted as published. \ifqa Statements from the question period
are quoted from a transcript of the webcast checked against the video, with
clip and time.\else Statements from the question period are not used here:
the transcript of the webcast has not yet been checked against the
video.\fi

\begin{small}
\begin{longtable}{>{\raggedright\arraybackslash}p{2.2cm}>{\raggedright\arraybackslash}p{4.4cm}>{\raggedright\arraybackslash}p{4.4cm}>{\raggedright\arraybackslash}p{3.6cm}}
\toprule
Topic & October~6 & \emph{Daybreak}, October~9 & The record \\
\midrule
\endhead
The ``8'' & ``We go from 7700 [700] during the good months to 8. That's 8 passage, right? So that's 4 cyclists, right? Including me, right?'' & ``it basically was eight during all of the hours. So, it was the sum of all the hours during the rush hour. So, that's six hours'' \cliptime{07:03} & Sum over six peak hours, one intersection (Grand), one day (January~29, 2025), both directions \pg{62}, \pg{34} \\
\addlinespace
Winter level & ``in winter, the volume of traffic diminishes by 99\%'' & ``less than one an hour'' \cliptime{03:11}; ``from 40 a day to 30 a day during December to March'' \cliptime{07:34} & Kensington counter: 40.4 a day December--March, 747 a day in June 2026 ($-95$\,\%) \\
\addlinespace
Data source & ``if we looked at the official data provided by the borough and we have counters everywhere, we have more counters than before. So we know exactly how many people use it, when and how'' & ``the source of truth is to look at the open data website that we have, which is connected directly to Eco-Counter'' \cliptime{07:19} & The 8 and the 700 are Miovision counts from the borough's slides; the Kensington Eco-Counter is the borough's and is not on the open data portal (section~2.1) \\
\addlinespace
Years of data & ``now having data for 3 years in a row'' & ``we have the data for many years to confirm that'' \cliptime{03:33} & One continuous winter (2025--26); single-day counts before \\
\addlinespace
Test or permanent & ``So this is what we're going to propose. It's going to be a test. We're going to see the data this year, and hopefully try to make it even better after.'' & ``So for now, we're looking at it in a way that will make it permanent'' \cliptime{12:22} & ``mesure transitoire pour l'hiver 2026-2027'' (August~20) \pg{4} \\
\addlinespace
Meetings & ``more than 30 meetings on this topic with the services, with professionals'' & ``more than 40 meetings and work sessions'' \cliptime{02:37} & Three meetings named in the release (section~2.5) \\
\addlinespace
Doors & ``door knocking on 1000s of doors'' & ``I've knocked on 4,000 doors'' \cliptime{01:53} & No record \\
\addlinespace
Parking spaces & No number in the address or on the statement card & ``at least 50 to 60 parking spots'' \cliptime{01:28} & Not in the release (section~2.2) \\
\addlinespace
Steps and timing & Card: ``several steps remain … borough council and City Council … residents will see changes before Christmas 2026'' & Borough vote, Conseil municipal vote, back to the borough; ``it's very tight''; ``implement those changes before Christmas'' \cliptime{08:19} & Consistent with each other and with the 2024 precedent (three votes) \\
\ifqa
\addlinespace
The ``8'' (question period) & Clip~04 (French): ``des journées où il y a huit personnes qui l'utilisent''; ``c'est huit comptages pour quatre cyclistes''. Clip~05: ``eight trips in January … Per hour. Heure de pointe.'' & As above: a six-hour sum \cliptime{07:03} & Six-hour sum \pg{62} \\
\addlinespace
The study & Clip~03 (French): ``on a reçu le dernier rapport qui nous permettait de prendre une décision en consensus la semaine passée''. Clip~05: ``we have the study … an internal working document'' & ``the final study by the services and the experts still need to be submitted to us'' \cliptime{00:38} & Final note due October~2 \pg{12}; withheld \\
\addlinespace
Consultation & Clip~05: ``we will not have another moment to consult now before the implementation'' & ``we're all about looking at the data, looking at the impact, consulting the population'' \cliptime{11:58} & -- \\
\addlinespace
Test & Clip~04 (French): ``on va le tester cette année'' & ``for now… permanent'' \cliptime{12:22} & As above \\
\fi
\bottomrule
\end{longtable}
\end{small}

\section{What this check does not establish}

It compares statements with records. It does not establish why any
statement was made, and it draws no conclusion about intent. Counter data have limits: an outage on May~29--31, 2026; the Kensington
westbound channel undercounting from about March~1, 2026; the open-data
detector at Girouard counting one direction, with gaps. The final EXP note
and the City's analysis, which may contain figures not in the release, are
withheld.

\section{Sources}

\begin{small}
\begin{longtable}{>{\raggedright\arraybackslash}p{4.2cm}>{\raggedright\arraybackslash}p{10.6cm}}
\toprule
Source & Where \\
\midrule
\endhead
Interview & CBC \emph{Daybreak}, October~9, 2026: \url{https://www.cbc.ca/listen/live-radio/1-15-daybreak-montreal/clip/16243474-terrebonne-bike-path-cdn-shut-winter}; transcript: \doc{2026-cbc-daybreak-interview/derived/2026-10-09_CBC_Daybreak_Teodoresco_transcript_en.md}{this repository} \\
Request 2026-581 & Borough response of October~9, 2026, bundle 01: \href{\pagesurl/2026-581/raw/01_emails_mandate_counts.pdf}{PDF}; summaries and counter tables in the \href{\arurl/tree/main/2026-581}{access-requests repository} \\
Requests 2025-697, 2026-300 & Kensington counter, hourly by direction: \href{\arurl/tree/main/2026-300}{2026-300} \\
City open data & \url{https://donnees.montreal.ca/dataset/cyclistes}; \url{https://donnees.montreal.ca/dataset/comptage-vehicules-pietons}; \url{https://donnees.montreal.ca/dataset/reseau-arteriel-administratif} (downloaded October~16, 2026) \\
EXP 2023 & \doc{2023-exp-mobility-study/raw/EXP_Etude_de_mobilite_rue_de_Terrebonne_2023.pdf}{Étude de mobilité -- rue de Terrebonne} (pp.~13, 41, 46, 54, 100) \\
Council records & \doc{2024-project-approval/}{2024 project approval} (CA24 170018, CM24 0179); CA24 170131: \url{https://ville.montreal.qc.ca/documents/Adi_Public/CA_Cdn/CA_Cdn_PV_ORDI_2024-06-03_19h00_FR.pdf} (p.~11) \\
Budget & \url{https://ville.montreal.qc.ca/documents/Adi_Public/CA_Cdn/CA_Cdn_ODJ_LPP_EXTRA_2026-10-06_16h00_FR.pdf} (pp.~14, 31) \\
Executive committee & \url{https://montreal.ca/communiques-de-presse/la-mairesse-martinez-ferrada-devoile-son-comite-executif} \\
Vélorues & \url{https://www.quebec.ca/transports/municipalites/infrastructures-routieres/velorue}; \url{https://montreal.ca/articles/velorues-des-voies-cyclables-accessibles-aux-automobilistes-54165} \\
October~6, 2026 & Borough council webcast (address; question period); councillor's statement card \\
Press & \emph{Métro}, September~1, 2020; \emph{Montreal Gazette}, November~26, 2025 and October~7, 2026 \\
\bottomrule
\end{longtable}
\end{small}

\end{document}
